\documentclass[12pt]{article}

\usepackage{fontawesome}
\usepackage{hyperref}
\usepackage{xurl}
\usepackage[tagged, highstructure]{accessibility}
\usepackage{bookmark}
\usepackage{enumitem}

\hypersetup{
    colorlinks=false,
    pdfborder={0 0 0},
    pdftitle={Editor (vi)},
    pdfauthor={Adrianna Holden-Gouveia},
    pdfsubject={Introduction to Linux - Lab Assignment},
    pdflang={en-US},
    pdfkeywords={lab assignment, Editor (vi)}
}

% Configure PDF bookmarks for navigation
\bookmarksetup{
    numbered,
    open,
}

% Configure list spacing for better accessibility
\setlist{nosep}



\title{Editor (vi)}
\author{
        Adrianna Holden-Gouveia \\
        Website: \url{https://aholdengouveia.name}\\ 
        \date{\vspace{-5ex}}
        %Email: \href{mailto:admin@aholdengouveia.name}{admin@aholdengouveia.name} \\
%        \faLinkedin{: aholdengouveia} \\
        \faGithub {: aholdengouveia} \\
%        \faTwitter {: aholdengouveia} \\
        }

%S\date{\today}


\begin{document}    

\maketitle

\section*{Accessibility Notice}
This document is also available in HTML format at:

\url{https://aholdengouveia.name/IntroLinux/labs/envvar.html}

The HTML version provides enhanced accessibility features including keyboard navigation, screen reader support, responsive design, dark mode support, and high contrast options.


%\begin{abstract}
%This is a template for Linux Administration Lab work
%\end{abstract}
%\tableofcontents

\section*{Objectives:}
\begin{itemize}
    \item Learning vi (editor) (Actually vim)
You must be able to demonstrate creating a file with text, closing and saving the file, Deleting text, inserting additional text, modifying text.
\end{itemize}


\section*{Complete the following problems}

References, a video, a PowerPoint and some notes are available at my website
\url {https://www.aholdengouveia.name/IntroLinux/envvar.html}

Practice using \url{https://vim-adventures.com/} or the VI Tutor that's built in to VI  or \url{https://www.openvim.com/}

VI is a modal editor, and as such it doesn't work the same as other text editing program you might have used, please make sure you have gone through the learning module before attempting this lab if you've never used VI or VIM before. 

\subsection*{Please include}
    \begin{enumerate}
        \item Using vi, create a file named Introduction that contains the following text (don’t worry about indenting just type the text or learn to paste in the terminal): 
        \begin{enumerate}[label=---,leftmargin=2em]
        \item vi (pronounce: ``vee eye'', not ``six'', not ``vye'') is an editor. An editor is a program to edit files. Goodbye.

            Although other stories exist, the true one tells that vi was originally written by Bill Joy in 1976. Bill took the sources of ed and ex, two horrendous programs for Unix that try to enable a human being to edit files, and created Vi. A truly remarkable, and somewhat paradoxical, event. Read the interview with Bill Joy for a more accurate history of Vi.

            People got attached to Vi, and eventually it got included in System V. From there on history has covered its traces and now Vi has evolved in many different versions for many, many platforms. The basic concept of Vi, however, has not changed over the years.
    \end{enumerate}


        \item Close and save the file.
        \item Re-open the file and add your name and the date at the top.
        \item Append to the bottom of the file the following text: 
            \begin{enumerate}[label=---,leftmargin=2em]
                \item"It's a vi thing - you wouldn't understand."
                
                Yes, vi is a religion. More than that, it's a cult religion. It started out as a UNIX editor and has now become a Ubiquitous editor. Unlike, say, Microsoft Word, it can also give the illusion of magic. Watching a vi guru doing some heavy editing on a file, as her fingers fly over the keys and textual transformations sweep across the screen, one could believe that one is in the presence of supernatural powers. "Awesome", the audience murmur to themselves. Maybe. \\
            \end{enumerate}


    \end{enumerate}
  Be sure you know how to : edit files, insert text, delete text, locate text, save and exit


\section*{Deliverables}
 Two files, one should be a clearly labeled screenshot of your file open in VI, and the second file should be the actual file from the server. 
\end{document}